%! AP Statistics — Unit 2, Lesson 2.6 — Student Packet Cover Sheet
\documentclass[10pt]{article}
\usepackage{apstats-article}
\usepackage{apstats-boxes}
\usepackage{ltablex}
\keepXColumns

\begin{document}

% Full-bleed navy banner
\begin{tikzpicture}[remember picture, overlay]
  \fill[navy] (current page.north west)
    rectangle ([yshift=-0.9in]current page.north east);
\end{tikzpicture}
\vspace{-0.8in}

\begin{center}
  {\color{white}\textbf{\LARGE AP Statistics}}\\[4pt]
  {\color{skymid}\large\textbf{Unit 2: Exploring Two-Variable Data}}\\[6pt]
  {\color{white}\textbf{\large Lesson 2.6 \quad Linear Regression Models}}
\end{center}

\vspace{0.22in}
\namedateperiod
\vspace{0.18in}

% ── LEARNING TARGETS ─────────────────────────────────────────────────────────
\begin{learningtargetbox}
\small
\begin{enumerate}[leftmargin=1.5em, itemsep=5pt]
  \item \ldots{} write the least-squares regression equation using the names of the variables
        (not generic $x$ and $y$). \hfill\textcolor{slate}{\small(DAT-1.E)}
  \item \ldots{} interpret the \textbf{slope} of the regression line in context, including
        direction and units. \hfill\textcolor{slate}{\small(DAT-1.E)}
  \item \ldots{} interpret the \textbf{\textit{y}-intercept} of the regression line in context
        and state whether it is meaningful. \hfill\textcolor{slate}{\small(DAT-1.E)}
  \item \ldots{} use the regression equation to make \textbf{predictions} and distinguish
        \textbf{interpolation} from \textbf{extrapolation}. \hfill\textcolor{slate}{\small(DAT-1.F)}
\end{enumerate}
\end{learningtargetbox}

\vspace{0.14in}

% ── TABLE OF CONTENTS ────────────────────────────────────────────────────────
\begin{tocbox}
\small
\renewcommand{\arraystretch}{1.5}
\begin{tabularx}{\linewidth}{@{} c l X r @{}}
\rowcolor{sky}
\textbf{\#} & \textbf{Component} & \textbf{Description} & \textbf{Score} \\
\midrule
1 & Warm-Up       & Spiral review --- slope-intercept form and correlation                & \blank{1.2cm} \\
2 & Guided Notes  & Regression equation, slope and intercept interpretation, predictions  & \blank{1.2cm} \\
3 & Group Activity & Interpreting and applying regression equations in context            & \blank{1.2cm} \\
4 & Exit Ticket   & Interpret and apply a regression equation independently               & \blank{1.2cm} \\
5 & Homework      & Multi-context regression interpretation and prediction practice        & \blank{1.2cm} \\
\midrule
  & \multicolumn{2}{r}{\textbf{Total}} & \blank{1.2cm} \\
\end{tabularx}
\end{tocbox}

\vspace{0.14in}

% ── KEEP IN MIND ─────────────────────────────────────────────────────────────
\begin{remindbox}
\small
The regression line $\hat{y} = a + bx$ is the tool for predicting one quantitative variable
from another. Here are the rules for full credit on the AP exam.

\vspace{6pt}
\begin{tabularx}{\linewidth}{@{} >{\centering\arraybackslash\columncolor{goldbg}}p{0.06\linewidth}
                                    X @{}}
\textbf{\large!} &
\textbf{Use variable names, not $x$ and $y$.}\enspace
Write $\widehat{\text{price}} = 28.4 - 1.85(\text{age})$, not $\hat{y} = 28.4 - 1.85x$.
The AP exam requires context in every interpretation. \\[6pt]
\textbf{\large!} &
\textbf{Slope interpretation must include direction, magnitude, and context.}\enspace
``For each additional year of age, the predicted resale price decreases by approximately
\$1{,}850.'' All three pieces are required. \\[6pt]
\textbf{\large!} &
\textbf{Check whether $x = 0$ is meaningful before interpreting the intercept.}\enspace
If no individual could realistically have $x = 0$, state that the intercept ``is not
meaningful in context'' --- do not skip this judgment call. \\[6pt]
\textbf{\large!} &
\textbf{Flag extrapolation.}\enspace
A prediction for an $x$ value outside the data range is extrapolation. It is unreliable
because the linear pattern may not extend that far. Interpolation (within the data range)
is generally trustworthy.
\end{tabularx}
\end{remindbox}

\end{document}
